\chapter{Нерадиальный спектр вихря после исправления кинетической нормы}

\section{Почему прежний сертификат не переносится}

Единый условный след исправил коэффициенты эффективной поперечной энергии:
\begin{equation}
 A=C=\frac{128}{243},\qquad
 B=\frac{160}{243},\qquad
 G=\frac{2}{27}.
 \label{eq:v6-corrected-angular-coefficients}
\end{equation}
Изменились не только общий масштаб скалярной кинетики, но и отношение
$A/G$, входящее в фоновую калибровку. Поэтому спектр старого профиля
нельзя получить умножением прежних собственных значений на константу.

Методическая граница задана литературой. Конечные вихри нарушения
$SO(3)\to A_4$ построены в \path{arXiv:2607.12366}. Возможность
длинноволновой неустойчивости вихревой нити показана в
\path{arXiv:0712.3589}. Фоновая калибровка и требование отделить ровно
геометрические нулевые моды от чистых калибровочных направлений разобраны
в \path{arXiv:1602.09084}; современный самосопряжённый спектральный
анализ абелева вихря дан в \path{arXiv:2608.10608}.

\section{Исправленный оператор}

На заново решённом фоне $(K,a,b)$ рассмотрена энергия
\begin{equation}
 E_\perp=\int_{\mathbb R^2}\left[
 \frac A2|D_i\phi|^2+\frac B2|\partial_i b|^2
 +\frac G2F_{12}^2+V(|\phi|,b)\right]d^2x,
\end{equation}
где $\phi_0=a(r)e^{i\theta}$ и $A_{0\theta}=K(r)/r$.
Для возмущений добавлена фоновая калибровка
\begin{equation}
 \mathcal G=\partial_i\alpha_i+\frac AG(J\phi_0)\mathbin{\cdot}\eta,
 \qquad \frac AG=\frac{64}{9}.
\end{equation}
Проверяемый оператор действует на
$(\Re\eta,\Im\eta,\alpha_x,\alpha_y,\delta b)$ и нормирован временной
метрикой с весами $(A,A,G,G,B)$.

\section{Сгущение декартовой сетки}

Расчёт выполнен в квадрате радиуса $10$ на сетках
$24,32,40,48,56$. На последней сетке размерность оператора равна
$14580$, а среднеквадратичный остаток дискретного уравнения фона равен
\begin{equation}
 7.7186033\cdot10^{-4}.
\end{equation}
Восемь нижних собственных значений равны
\begin{equation}
\begin{split}
 &(0.01327643,\ 0.01331942,\ 0.01391266,\ 0.01391579,\\
 &\hspace{1.4cm}0.01897747,\ 0.01948457,\ 0.02431133,\ 0.02436721).
\end{split}
\end{equation}
Отрицательных значений на проверенных сетках нет. Однако это не является
сертификатом устойчивости: при переходе от сетки $48$ к сетке $56$
третья мода уменьшается в
\begin{equation}
 \frac{\lambda_3^{(56)}}{\lambda_3^{(48)}}
 =0.1105014\ldots,
\end{equation}
а все восемь проверенных мод входят в окрестность нуля $|\lambda|<0.03$.
Следовательно, прежнее выделение двух мягких переносов и положительного
жёсткого зазора не выдерживает сгущения сетки после исправления нормы.

Прямое дифференцирование фона также не решает вопрос: проекции двух
наивных переносных касательных на четыре нижние моды падают до
$3.23\cdot10^{-5}$. Это означает не отсутствие переносов, а
необходимость добавить к геометрическому переносу компенсирующее
калибровочное преобразование и построить точный ковариантный проектор.

\section{Статус результата}

\begin{nogocriterion}
Численная неотрицательность конечной сетки не считается устойчивостью,
если более двух собственных значений стремительно смягчаются, переносные
касательные не отождествлены ковариантно и положительный зазор не
сходится. До разрешения нулевого подпространства переход к полному
тензорному гессиану и к замкнутой хопфовой петле преждевременен.
\end{nogocriterion}

Нулевая гипотеза следующего гейта состоит в том, что множественное
смягчение является недостатком дискретной фоновой калибровки. Она будет
отвергнута, если после точного ковариантного проектирования останется
более двух физических нулевых мод либо возникнет отрицательный предел.

\begin{protocol}
\begin{enumerate}
 \item исправленный нерадиальный оператор: \textbf{построен};
 \item отрицательные моды на проверенных сетках: \textbf{не найдены};
 \item сходящийся положительный зазор: \textbf{не получен};
 \item число смягчающихся проверенных мод: \textbf{не менее восьми};
 \item ковариантное выделение переносов: \textbf{не выполнено};
 \item нерадиальная устойчивость: \textbf{не сертифицирована};
 \item следующий гейт: \textbf{ковариантное разрешение нулевых мод}.
\end{enumerate}
\end{protocol}

Машинный сертификат сохранён в
\path{s2t_v6_bosonic_defect_corrected_vortex_nonradial_stability_gate_results.json}.